\section*{第 \ref{ch:g10:reaction-progress-table} 章——反应进度表}

\begin{solution}{exo:g10:reaction-progress-table:1}
\ce{H2}：$4.0 - 2x$；\ce{O2}：$3.0 - x$；\ce{H2O}：$2x$。
氢气在 $x = 2.0$ 时耗尽，氧气在 $x = 3.0$ 时耗尽，故 $x_{\max} = \qty{2.0}{mol}$，
氢气为限量反应物。终态：\qty{0}{mol} \ce{H2}，
\qty{1.0}{mol} \ce{O2}，\qty{4.0}{mol} \ce{H2O}。
\end{solution}

\begin{solution}{exo:g10:reaction-progress-table:2}
碳将在 $x = 3.0$ 时耗尽，氧气将在 $x = 2.0$ 时耗尽：
$x_{\max} = \qty{2.0}{mol}$，氧气为限量反应物。终态：
剩余碳 \qty{1.0}{mol}，氧气耗尽，二氧化碳 \qty{2.0}{mol}。
\end{solution}

\begin{solution}{exo:g10:reaction-progress-table:3}
$x_{\max} = \qty{0.30}{mol}$（硫为限量反应物）。终态：
剩余铁 \qty{0.20}{mol}，硫耗尽，生成硫化亚铁 \qty{0.30}{mol}。
\end{solution}

\begin{solution}{exo:g10:reaction-progress-table:4}
氮气在 $x = 1.0$ 时耗尽，氢气在 $x = 2.4/3 = 0.80$ 时耗尽：
$x_{\max} = \qty{0.80}{mol}$。终态：\qty{0.20}{mol} \ce{N2}，
无 \ce{H2}，$2 \times 0.80 = \qty{1.6}{mol}$ \ce{NH3}。
\end{solution}

\begin{solution}{exo:g10:reaction-progress-table:5}
(b)：$4/2 = 2/1$。在 (a) 中一氧化碳为限量反应物，(c) 中亦然。
\end{solution}

\begin{solution}{exo:g10:reaction-progress-table:6}
$n(\ce{CH4}) = 32 / 16.0 = \qty{2.0}{mol}$，故 $x_{\max} = \qty{2.0}{mol}$
对于化学计量比混合物。氧气：$2 \times 2.0 = \qty{4.0}{mol}$，
即 $4.0 \times 32.0 = \qty{128}{g}$。水：
$\qty{4.0}{mol} \times \qty{18.0}{g/mol} = \qty{72}{g}$。
\end{solution}

\begin{solution}{exo:g10:reaction-progress-table:7}
$n(\ce{C3H8}) = 1000 / 44.0 = \qty{22.7}{mol}$；
二氧化碳 $3 \times 22.7 = \qty{68.2}{mol}$；$V = 68.2 \times 24.1 =
\qty{1.64e3}{L}$，约 \qty{1.6}{m^3}。
\end{solution}

\begin{solution}{exo:g10:reaction-progress-table:8}
当 $x = \qty{1.0}{mol}$ 时：\qty{1.0}{mol} \ce{CH4}，\qty{1.0}{mol}
\ce{O2}，\qty{1.0}{mol} \ce{CO2}，\qty{2.0}{mol} \ce{H2O}。氧气
在 $x = \qty{1.5}{mol}$ 时耗尽，反应在此停止，故更大的
反应进度永远不会达到。
\end{solution}

\begin{solution}{exo:g10:reaction-progress-table:9}
$n_0(\ce{Zn}) = 1.30 / 65.4 = \qty{0.0199}{mol}$；
$n_0(\ce{Cu^{2+}}) = 0.10 \times 0.1000 = \qty{0.0100}{mol}$。铜离子
为限量反应物：$x_{\max} = \qty{0.0100}{mol}$。生成的铜：
$0.0100 \times 63.5 = \qty{0.635}{g}$；剩余的锌：
$(0.0199 - 0.0100) \times 65.4 = \qty{0.65}{g}$。无铜离子剩余：
蓝色已消失。
\end{solution}

\begin{solution}{exo:g10:reaction-progress-table:10}
$n_0(\ce{Mg}) = 0.48 / 24.3 = \qty{0.0198}{mol}$，镁将在
$x = 0.0198$ 时耗尽；$n_0(\ce{H+}) = \qty{0.020}{mol}$，氢离子将在
$x = 0.010$ 时耗尽。酸为限量反应物：$x_{\max} = \qty{0.010}{mol}$，故
生成氢气 \qty{0.010}{mol}，$0.010 \times 24.1 = \qty{0.24}{L}$。
\end{solution}

\begin{solution}{exo:g10:reaction-progress-table:11}
化学计量比混合物中含 \qty{0.30}{mol} 硫，需要
\qty{0.30}{mol} 铁；铁多 \qty{0.20}{mol}，即
$0.20 / 0.30 \approx \qty{67}{\percent}$ 过量。
\end{solution}

\begin{solution}{exo:g10:reaction-progress-table:12}
物质的量相等：$m(\ce{Fe})/55.8 = m(\ce{S})/32.1$，其中
$m(\ce{Fe}) + m(\ce{S}) = \qty{10.0}{g}$。于是
$m(\ce{Fe}) = 10.0 \times 55.8 / (55.8 + 32.1) = \qty{6.35}{g}$，且
$m(\ce{S}) = \qty{3.65}{g}$。
\end{solution}

\begin{solution}{exo:g10:reaction-progress-table:13}
\begin{enumerate}
  \item $n_0(\ce{CaCO3}) = 10.0 / 100.1 = \qty{0.0999}{mol}$，唯一的
        反应物：$x_{\max} = \qty{0.0999}{mol}$；生成
        \qty{0.0999}{mol} 的 \ce{CaO} 和 \ce{CO2}，即
        $0.0999 \times 24.1 = \qty{2.41}{L}$ 的二氧化碳。
  \item $n_0(\ce{CaO}) = \qty{0.0999}{mol}$，$n_0(\ce{H2O}) = 1.8 / 18.0
        = \qty{0.100}{mol}$：生石灰（恰好）为限量反应物；
        混合物接近化学计量比。\ce{Ca(OH)2}：
        $M = 40.1 + 2 \times (16.0 + 1.0) = \qty{74.1}{g/mol}$，质量
        为 $0.0999 \times 74.1 = \qty{7.40}{g}$。
\end{enumerate}
\end{solution}

\begin{solution}{exo:g10:reaction-progress-table:14}
$n_0(\ce{H+}) = \qty{0.050}{mol}$（每只烧瓶中）；氢离子将在
$x = \qty{0.025}{mol}$ 时耗尽。镁：0.15、0.30、0.45、0.60、0.75、
\qty{0.90}{g} 分别得到 0.0062、0.0123、0.0185、0.0247、0.0309、
\qty{0.0370}{mol}。前四只烧瓶中镁为限量反应物，
氢气体积为 $n(\ce{Mg}) \times 24.1$：0.15、0.30、0.45 和
\qty{0.60}{L}。后两只烧瓶中酸为限量反应物：$0.025 \times 24.1
= \qty{0.60}{L}$。图像为过原点的直线，直至
约 \qty{0.61}{g} 镁，然后是在
\qty{0.60}{L} 处的水平平台。
\end{solution}

\begin{solution}{exo:g10:reaction-progress-table:15}
$n_0(\text{乙醇}) = 4.6 / 46.0 = \qty{0.100}{mol}$（将在
$x = 0.100$ 时耗尽）；$n_0(\ce{O2}) = 4.8 / 24.1 = \qty{0.199}{mol}$（将在
$x = 0.199/3 = 0.0664$ 时耗尽）。氧气为限量反应物，
$x_{\max} = \qty{0.0664}{mol}$。剩余乙醇：$0.100 - 0.0664 =
\qty{0.0336}{mol}$，即 $0.0336 \times 46.0 = \qty{1.5}{g}$。
二氧化碳：$2 \times 0.0664 = \qty{0.133}{mol}$，即
$0.133 \times 24.1 = \qty{3.2}{L}$。
\end{solution}

\begin{solution}{pb:g10:reaction-progress-table:1}
\textbf{1.} 左边：2 个 Na，6 个 N。右边：2 个 Na，$3 \times 2 = 6$ 个 N。已配平。

\textbf{2.} 左边：10 个 Na，2 个 K，2 个 N，6 个 O。右边：K：2；Na：$5 \times 2 =
10$；O：$1 + 5 = 6$；N：2。已配平。

\textbf{3.} 参见元素周期表章节（\cref{ch:g9:periodic-table-first-look}）：
钠与水剧烈反应，生成具有腐蚀性的溶液和
氢气。安全气囊破裂后，其残留物会接触皮肤、
眼睛以及呼出的潮湿空气。

\textbf{4.} $M(\ce{NaN3}) = 23.0 + 3 \times 14.0 = \qty{65.0}{g/mol}$；
$M(\ce{KNO3}) = 39.1 + 14.0 + 3 \times 16.0 = \qty{101.1}{g/mol}$.

\textbf{5.} \ce{NaN3}：$1.00 - 2x$；\ce{Na}：$2x$；\ce{N2}：$3x$。

\textbf{6.} 当 $1.00 - 2x = 0$ 时，$x_{\max} = \qty{0.500}{mol}$；叠氮
化钠为唯一的反应物，故为限量反应物。

\textbf{7.} \ce{Na}：$2 \times 0.500 = \qty{1.00}{mol}$；\ce{N2}：
$3 \times 0.500 = \qty{1.50}{mol}$。

\textbf{8.} $1.50 \times 24.1 = \qty{36.2}{L}$。

\textbf{9.} \ce{Na}：$1.00 - 10x$；\ce{KNO3}：$n - 2x$；\ce{K2O}：$x$；
\ce{Na2O}：$5x$；\ce{N2}：$x$。

\textbf{10.} $1.00/10 = n/2$，故 $n = \qty{0.200}{mol}$，即
$0.200 \times 101.1 = \qty{20.2}{g}$。

\textbf{11.} $x_{\max} = \qty{0.100}{mol}$：钠、硝酸钾均耗尽，生成
\qty{0.100}{mol} \ce{K2O}、\qty{0.500}{mol} \ce{Na2O}、
\qty{0.100}{mol} \ce{N2}。

\textbf{12.} 钠将在 $x = 0.100$ 时耗尽，硝酸钾将在 $x = 0.150/2 = 0.075$ 时耗尽：
硝酸钾为限量反应物，$x_{\max} =
\qty{0.075}{mol}$，袋中将残留 $1.00 - 10 \times 0.075 = \qty{0.25}{mol}$
的金属钠，遇湿气即会反应。

\textbf{13.} $1.50 + 0.100 = \qty{1.60}{mol}$，即每摩尔叠氮
化钠对应的氮气量。

\textbf{14.} $n = 60 / 24.1 = \qty{2.49}{mol}$。

\textbf{15.} $2.49 / 1.60 = \qty{1.56}{mol}$ 的叠氮化钠。

\textbf{16.} $1.56 \times 65.0 = \qty{101}{g}$。

\textbf{17.} $0.200 \times 1.56 = \qty{0.311}{mol}$，即
$0.311 \times 101.1 = \qty{31.5}{g}$ 的硝酸钾。

\textbf{18.} $2.49 / 1.50 = \qty{1.66}{mol}$，即
$1.66 \times 65.0 = \qty{108}{g}$：第二个反应可节省约
\qty{7}{g} 的叠氮化钠。

\textbf{19.} 完全反应仅在限量反应物耗尽时才会停止，
此处叠氮化钠为唯一的反应物，故终态下
其物质的量为 $1.00 - 2x_{\max} = 0$。袋中不再残留有毒物质。

\textbf{20.} 约 \qty{101}{g} 叠氮化钠可充满一个 \qty{60}{L}
的安全气囊。
\end{solution}
